\section{Cycle activity and phase-uniform determinant certificates}
\label{sec:tpc67-cycle-activity}

Retain the matching-normalized core \(I+E_\phi\) and the directed graph
\(\mathcal D_\phi\) from \cref{sec:tpc67-forced-core}.  Absolute values of
matrices in this section are entrywise.  The point is to separate two
questions:
\begin{enumerate}[label=(\roman*)]
\item the exact total mass of all nonidentity determinant terms; and
\item coefficient-blind upper bounds for that mass.
\end{enumerate}
The first has the sharp threshold one.  A convenient bound through the sum
of individual cycle weights has the smaller sufficient threshold
\(\log 2\), which is not sharp.

\subsection{The exact permanent activity}

\begin{definition}[Matching activity]
\label{def:tpc67-matching-activity}
The phase-blind activity of the matching \(\phi\) is
\begin{equation}
 \mathcal A_\phi
 :=
 \per(I+|E_\phi|)-1.
 \label{eq:tpc67-activity-permanent}
\end{equation}
\end{definition}

\begin{proposition}[Exact activity identity]
\label{prop:tpc67-exact-activity}
The activity is exactly the total absolute weight of all nonempty
vertex-disjoint directed-cycle families:
\begin{equation}
 \boxed{
 \mathcal A_\phi
 =
 \sum_{\mathcal F\in\mathfrak F_\phi}
 |w_\phi(\mathcal F)|.}
 \label{eq:tpc67-activity-family-sum}
\end{equation}
\end{proposition}

\begin{proof}
Expand the permanent of \(I+|E_\phi|\).  Its identity term is one.
The correspondence in \cref{thm:tpc67-cycle-family-expansion} applies
again, now without permutation signs.  The term attached to
\(\mathcal F\) is precisely
\[
 \prod_{C\in\mathcal F}\prod_{(i,j)\in C}|(E_\phi)_{i,j}|
 =
 |w_\phi(\mathcal F)|.
\]
Subtracting the identity term proves
\eqref{eq:tpc67-activity-family-sum}.
\end{proof}

\begin{theorem}[Sharp phase-uniform activity certificate]
\label{thm:tpc67-activity-certificate}
Let
\[
 P_{\mathrm{for}}
 :=
 \prod_{e\in\mathcal P_{\mathrm{for}}}|t_e|,
 \qquad
 P_\phi:=\prod_{i\in V}|p_i|.
\]
Then
\begin{equation}
 \boxed{
 |\det T|
 \ge
 P_{\mathrm{for}}P_\phi
 (1-\mathcal A_\phi)_+.}
 \label{eq:tpc67-activity-determinant-bound}
\end{equation}
In particular, \(\mathcal A_\phi<1\) implies that the literal minor is
nonsingular, uniformly in all phases having the prescribed entry
magnitudes.
\end{theorem}

\begin{proof}
By \eqref{eq:tpc67-exact-cycle-family-expansion}, the identity term of
\(\det(I+E_\phi)\) is one and the sum of the absolute values of all other
terms is \(\mathcal A_\phi\).  The reverse triangle inequality gives
\[
 |\det(I+E_\phi)|
 \ge (1-\mathcal A_\phi)_+.
\]
Multiplication by the exact pivot factors in
\eqref{eq:tpc67-full-core-normalization} proves the result.
\end{proof}

\begin{proposition}[Why the threshold one is sharp]
\label{prop:tpc67-activity-threshold-sharp}
Neither the threshold \(1\) nor the coefficient in
\((1-\mathcal A_\phi)_+\) can be improved using only the total activity.
Indeed, for
\begin{equation}
 I+E=
 \begin{pmatrix}
 1&x\\
 y&1
 \end{pmatrix}
 \label{eq:tpc67-two-cycle-model}
\end{equation}
one has
\[
 \mathcal A_\phi=|xy|,
 \qquad
 \det(I+E)=1-xy.
\]
Taking \(x,y>0\) gives equality in
\eqref{eq:tpc67-activity-determinant-bound}; at \(xy=1\), the determinant
vanishes and \(\mathcal A_\phi=1\).  On the other hand,
\(\mathcal A_\phi\ge1\) does not imply singularity: taking \(xy\le-1\) gives
\(|\det(I+E)|=1+\mathcal A_\phi\).
\end{proposition}

\subsection{From single-cycle mass to family activity}

Let \(\mathscr C_\phi\) be the set of simple directed cycles of
\(\mathcal D_\phi\), and define the single-cycle mass
\begin{equation}
 \Xi_\phi
 :=
 \sum_{C\in\mathscr C_\phi}
 |w_\phi(C)|.
 \label{eq:tpc67-single-cycle-mass}
\end{equation}

\begin{theorem}[Exponential domination of disjoint-cycle activity]
\label{thm:tpc67-exponential-activity}
One has
\begin{equation}
 \boxed{
 \mathcal A_\phi
 \le
 \exp(\Xi_\phi)-1.}
 \label{eq:tpc67-activity-exponential-bound}
\end{equation}
Consequently, if
\begin{equation}
 \Xi_\phi<\log 2,
 \label{eq:tpc67-logtwo-condition}
\end{equation}
then
\begin{equation}
 \boxed{
 |\det T|
 \ge
 P_{\mathrm{for}}P_\phi
 \bigl(2-e^{\Xi_\phi}\bigr)>0.}
 \label{eq:tpc67-logtwo-determinant-bound}
\end{equation}
\end{theorem}

\begin{proof}
Expand the exponential:
\[
 e^{\Xi_\phi}
 =
 \sum_{k=0}^{\infty}
 \frac1{k!}
 \left(
 \sum_{C\in\mathscr C_\phi}|w_\phi(C)|
 \right)^k.
\]
For a family of \(k\) distinct vertex-disjoint cycles, its weight appears
exactly \(k!\) times in the \(k\)-th power, once for every ordering of its
cycles, and hence with coefficient one after division by \(k!\).  The
exponential also contains nonnegative terms from repeated or intersecting
cycles.  It therefore majorizes one plus the family sum in
\eqref{eq:tpc67-activity-family-sum}, which proves
\eqref{eq:tpc67-activity-exponential-bound}.  If
\(\Xi_\phi<\log2\), then
\(\mathcal A_\phi\le e^{\Xi_\phi}-1<1\); substituting this estimate into
\eqref{eq:tpc67-activity-determinant-bound} gives
\eqref{eq:tpc67-logtwo-determinant-bound}.
\end{proof}

\begin{remark}[The \(\log2\) threshold is sufficient, not sharp]
\label{rem:tpc67-logtwo-not-sharp}
The sharp phase-blind threshold applies to
\(\mathcal A_\phi\), not to its exponential majorant.  In the two-cycle
model \eqref{eq:tpc67-two-cycle-model}, put
\[
 xy=q\in(\log2,1).
\]
There is only one simple cycle, so
\[
 \Xi_\phi=\mathcal A_\phi=q,
 \qquad
 \det(I+E)=1-q>0.
\]
Thus the exact activity certificate succeeds although
\(\Xi_\phi<\log2\) fails.  The loss comes solely from allowing arbitrary
ordered tuples of repeated and intersecting cycles in the exponential.
\end{remark}

\subsection{A degree--entry-size bound}

Put
\begin{align}
 q_\phi
 &:=
 |\{(i,j):i\ne j,\ (E_\phi)_{i,j}\ne0\}|,
 \nonumber\\
 \Delta_\phi^+
 &:=
 \max_{i\in V}
 |\{j\ne i:(E_\phi)_{i,j}\ne0\}|,
 \nonumber\\
 \eta_\phi
 &:=
 \max_{i\ne j}|(E_\phi)_{i,j}|,
 \label{eq:tpc67-degree-entry-parameters}
\end{align}
with \(\eta_\phi=0\) when \(q_\phi=0\), and write \(n=|V|\).

\begin{proposition}[Cycle mass from out-degree and pivot ratios]
\label{prop:tpc67-degree-entry-cycle-bound}
For every matching-normalized core,
\begin{equation}
 \boxed{
 \Xi_\phi
 \le
 \sum_{k=2}^{n}
 \frac{q_\phi}{k}
 (\Delta_\phi^+)^{k-2}\eta_\phi^k.}
 \label{eq:tpc67-finite-degree-entry-bound}
\end{equation}
If \(\Delta_\phi^+\eta_\phi<1\), then
\begin{align}
 \Xi_\phi
 &\le
 \frac{q_\phi\eta_\phi^2}
      {2(1-\Delta_\phi^+\eta_\phi)}
 \label{eq:tpc67-geometric-cycle-bound}\\
 &\le
 \frac{n\Delta_\phi^+\eta_\phi^2}
      {2(1-\Delta_\phi^+\eta_\phi)}.
 \label{eq:tpc67-coarse-geometric-cycle-bound}
\end{align}
Consequently,
\begin{equation}
 \mathcal A_\phi
 \le
 \exp\left\{
 \frac{q_\phi\eta_\phi^2}
      {2(1-\Delta_\phi^+\eta_\phi)}
 \right\}-1,
 \label{eq:tpc67-degree-entry-activity-bound}
\end{equation}
and the explicit sufficient condition
\begin{equation}
 \frac{q_\phi\eta_\phi^2}
      {2(1-\Delta_\phi^+\eta_\phi)}
 <\log2
 \label{eq:tpc67-degree-entry-logtwo}
\end{equation}
implies the positive determinant margin
\eqref{eq:tpc67-logtwo-determinant-bound}.
\end{proposition}

\begin{proof}
For a directed simple cycle of length \(k\), choose one of its \(k\)
oriented edges as the starting edge.  There are at most \(q_\phi\) choices
for that edge and at most
\((\Delta_\phi^+)^{k-2}\) choices for the next \(k-2\) edges; the closing
edge, if present, is then fixed.  Since each cycle was counted \(k\) times,
the number of length-\(k\) cycles is at most
\[
 \frac{q_\phi}{k}(\Delta_\phi^+)^{k-2}.
\]
Every such cycle has absolute weight at most \(\eta_\phi^k\), proving
\eqref{eq:tpc67-finite-degree-entry-bound}.

When \(\Delta_\phi^+\eta_\phi<1\), use \(1/k\le1/2\) for \(k\ge2\) and
sum the resulting geometric series.  This gives
\eqref{eq:tpc67-geometric-cycle-bound}.  The inequality
\(q_\phi\le n\Delta_\phi^+\) gives
\eqref{eq:tpc67-coarse-geometric-cycle-bound}.  Finally apply
\cref{thm:tpc67-exponential-activity}.
\end{proof}

\begin{remark}[What the combinatorial bound does not supply]
\label{rem:tpc67-cycle-activity-boundary}
The entry parameter is the pivot ratio
\[
 \eta_\phi
 =
 \max_{i\ne j}
 \frac{|t_{\phi(i),j}|}{|t_{\phi(j),j}|},
\]
not merely the largest unnormalized collision coefficient.  Moreover,
TPC--66's reduced pair-codegree estimate does not by itself bound
\(q_\phi\), \(\Delta_\phi^+\), or \(\eta_\phi\) for an actual full
\((\gamma,s,j)\)-minor \citep{WangTPC66}.  The present activity estimates
are therefore L0 certificates.  They become L1 only after the literal
minor, matching, pivot ratios, and full decorated collision core have been
identified.  No polynomial physical determinant bound, parity improvement,
or prime-pair conclusion is asserted here.
\end{remark}
